\chapter{Heatwaves}
\index{Heatwaves}

\section*{Definition and Overview}
A heatwave is a prolonged period of abnormally hot weather, defined in many countries as three or more days of unusually high maximum and minimum temperatures for the region. Heatwaves are among the most dangerous natural hazards in many countries, causing more deaths than floods, bushfires, and storms combined. They place enormous strain on the body, particularly for older people, young children, and those with chronic illness, and they also strain infrastructure such as power, water, and transport. Understanding and acting on heatwave warnings saves lives.

\section*{Epidemiology}
Heatwaves kill more people than any other natural hazard, with major events such as the 2009 Victorian heatwave contributing to more than 370 deaths. During heatwaves, hospital admissions and emergency presentations for heat illness, kidney failure, heart disease, and respiratory problems rise sharply. Older adults, particularly those living alone without air-conditioning, account for the majority of deaths. Climate change is making heatwaves more frequent, longer, and more intense, and heat-related mortality is projected to increase substantially in coming decades.

\section*{Aetiology and Pathophysiology}
\begin{itemize}
    \item \textbf{Heat exposure:} sustained high daytime and night-time temperatures prevent the body from cooling, especially overnight
    \item \textbf{Failed cooling:} high humidity and lack of airflow reduce the effectiveness of sweating
    \item \textbf{Vulnerable physiology:} older adults have reduced thirst, sweat responses, and cardiovascular reserve; infants cannot regulate temperature well
    \item \textbf{Chronic disease:} heart, kidney, lung, and metabolic disease are worsened by heat stress
    \item \textbf{Medications:} diuretics, blood pressure medicines, and some psychiatric drugs impair heat tolerance
    \item \textbf{Dehydration:} fluid loss through sweating leads to reduced blood volume and strain on the heart
    \item \textbf{Urban heat island:} cities are hotter than surrounding areas, increasing night-time temperatures
\end{itemize}

\section*{Clinical Features}
\begin{itemize}
    \item Heat exhaustion: heavy sweating, fatigue, dizziness, headache, nausea, and muscle cramps
    \item Heat stroke: very high temperature, confusion, seizures, and collapse, requiring emergency care
    \item Worsening of heart disease, with chest pain, breathlessness, and heart failure decompensation
    \item Kidney injury from dehydration
    \item Worsening of asthma and COPD
    \item Disturbed sleep, irritability, and reduced function in the community
    \item Death, often in older people found collapsed at home without cooling
\end{itemize}

\section*{Differential Diagnosis}
\begin{itemize}
    \item Heat illness must be distinguished from infection and fever
    \item Cardiac events, stroke, and other emergencies can mimic or coexist with heat illness
    \item Dehydration and gastroenteritis cause similar weakness and collapse
    \item Medication side effects may be mistaken for heat symptoms
    \item The environmental context of a heatwave makes heat-related illness the leading consideration
\end{itemize}

\section*{When to Seek Medical Care}
Seek medical advice for anyone with symptoms of heat illness that do not resolve quickly with cooling and fluids, and for vulnerable people who become unwell during a heatwave.

\textbf{Contact your local emergency services for:}
\begin{itemize}
    \item Confusion, agitation, seizures, or unconsciousness in the heat
    \item Body temperature of 40 degrees Celsius or higher
    \item Collapse, severe weakness, or inability to stand
    \item Chest pain, severe breathlessness, or signs of stroke
    \item A baby, older person, or person with chronic illness who is severely unwell
    \item Anyone who does not improve with cooling within 30 minutes
\end{itemize}

\section*{Diagnosis and Investigations}
\begin{itemize}
    \item \textbf{Heatwave warnings:} the Bureau of Meteorology issues heatwave forecasts and warnings that trigger public health alerts
    \item \textbf{Community assessment:} public health agencies monitor emergency presentations and deaths during heatwaves
    \item \textbf{Clinical assessment:} temperature, hydration, and organ function in affected people
    \item \textbf{Blood tests:} electrolytes, kidney function, and cardiac markers where clinically indicated
    \item \textbf{Home and environmental factors:} assessment of cooling availability and living conditions for vulnerable people
\end{itemize}

\section*{Management}
\begin{itemize}
    \item \textbf{Public health response:} heatwave warnings, cool refuges, extended opening hours, and targeted outreach to vulnerable groups
    \item \textbf{Personal response:} stay in cool, air-conditioned places; use fans; take cool showers; and limit activity
    \item \textbf{Hydration:} drink water regularly, even without thirst
    \item \textbf{Checking on others:} visit or phone older relatives, neighbours, and people living alone
    \item \textbf{Keeping homes cool:} shut curtains during the day, open windows at night, and use awnings and shade
    \item \textbf{Medical care:} treatment of heat illness and exacerbations of chronic disease
    \item \textbf{Infrastructure support:} power and water providers prepare for increased demand
\end{itemize}

\section*{Complications}
\begin{itemize}
    \item Heat stroke, with risk of permanent organ damage and death
    \item Increased heart attacks, strokes, and heart failure during heatwaves
    \item Kidney failure from dehydration
    \item Worsening of respiratory, diabetic, and neurological conditions
    \item Falls and injuries in older people
    \item Deaths among the most vulnerable, especially older people living alone
    \item Mental health effects, including anxiety, sleep disturbance, and isolation
\end{itemize}

\section*{Prognosis and Outcomes}
Most people manage heatwaves with sensible precautions and recover fully from mild heat illness. The key to good outcomes is preparation: acting on warnings before the heat arrives, cooling homes overnight, and checking on vulnerable people. Deaths during heatwaves are almost always preventable, and communities with heat plans, warning systems, and outreach to at-risk groups consistently fare better. Early cooling and treatment of heat illness prevent progression to life-threatening heat stroke.

\section*{Prevention and Self-Care}
\begin{itemize}
    \item Follow heatwave forecasts and warnings and act before the heat arrives
    \item Identify air-conditioned places, including shopping centres and libraries, and use them
    \item Keep homes cool: shade windows, use awnings, and open windows for night ventilation
    \item Drink water regularly and avoid alcohol and caffeinated drinks
    \item Plan outdoor activity for the coolest part of the day and rest often
    \item Never leave children or pets in parked cars
    \item Check on older relatives, neighbours, and people with chronic illness daily
    \item Prepare a heatwave plan for yourself and your household
\end{itemize}

\section*{Sources and Further Reading}
The following reputable sources provide further information for healthcare students and patients seeking reliable, current advice about heatwaves.

\begin{itemize}
    \item Australian Government Bureau of Meteorology. \textit{Heatwave warnings.} http://www.bom.gov.au/
    \item Healthdirect. \textit{Staying safe in a heatwave.} https://www.healthdirect.gov.au/heatwave
    \item Better Health Channel. \textit{Beat the heat: heatwave safety.} https://www.betterhealth.vic.gov.au/
    \item State and territory health departments. \textit{Heat health plans and alerts.} https://www.health.gov.au/
    \item World Health Organization. \textit{Heat and health: fact sheet.} https://www.who.int/
    \item Climate Council. \textit{Heatwaves and climate change in Australia.} https://www.climatecouncil.org.au/
\end{itemize}
